\documentclass[12pt]{article}
%%%%%%%%%%%%%%%%%%%%%%%%%%%%%%%%%%%%%%
%    Notes
%%%%%%%%%%%%%%%%%%%%%%%%%%%%%%%%%%%%%%
% This is a template we have designed to facilitate the reporting of whole genome sequencing-based results for Mycobacterium tuberculosis diagnosis, phenotyping, and epidemiologcal clustering. A manuscript describing how we arrived at this template is being submitted to PeerJ in October, 2017. We will update this document with a link to the article – "Evidence-Based Design and Evaluation of a Whole Genome Sequencing Clinical Report for the Reference Microbiology Laboratory" – when it comes online.

% A working example of using R code to automatically populate this report with your own data is available at: https://github.com/amcrisan/TB-WGS-MicroReport. Different laboratories will handle data in different ways, so your experience in implementing the template may vary. 

% If you are interested in using this template in your own laboratory, we recommend working with your staff bioinformatician to implement the automatic populating of the fields from your WGS analysis pipeline.

% This template is made freely available to the community and we encourage you to modify it to suit your own laboratory's needs, and to use it for pathogens other than M. tuberculosis. If you publish your pipeline and include a description of the report, we would appreciate a citation to the PeerJ paper (see above)

% Contact :
% For template-related questions, please contact Ana Crisan at acrisan@cs.ubc.ca
% For larger project-related questions, please contact Dr. Jennifer Gardy at jennifer.gardy@bccdc.ca


% Version 1.0 : 25 SEPT. 2017
% Version 1.1 : 1 OCT. 2017

%%%%%%%%%%%%%%%%%%%%%%%%%%%%%%%%%%%%%%
%    Version Change Notes
%%%%%%%%%%%%%%%%%%%%%%%%%%%%%%%%%%%%%%

%-----------------------
%  Version 1.0
%
% Initial release based upon PEERJ paper study and results

%-----------------------
%  Version 1.1
%
% Updates from consortium update, added the following:
% Header content section: 
%- delete ‘WHOLE’ so it just says ‘GENOME SEQUENCING REPORT’ - the report may be used to share the results of targeted sequencing - just amplicons containing resistance genes - so easiest to ditch the ‘WHOLE’

%2. The table with the patient/sample info has a few new fields and some changes: 
%- a place for a barcode from LIMS systems has been added
%- fields for sample source, sample ID, and material that was sequenced have been added

%3. Mycobacterium tuberculosis has been italicized throughout

%4. In the Organism box, we have added the lineage to which the specimen belongs 

%5. In the Drug Susceptibility box:
%- the explanatory paragraph has been modified slightly
%- 'Prediction’ has been changed to ‘Interpretation’
%- ‘Sensitive’ has been changed to ‘Susceptible’
%- a missing ‘x’ in moxifloxacin was added
%- 'No resistance mutation detected’ has been shortened to 'No mutation detected’ 

%6. On the second page: 
%- an Assay Details and Disclaimer section have been added
%- I had to monkey with the spacing to get everything back onto one page



%%%%%%%%%%%%%%%%%%%%%%%%%%%%%%%%%%%%%%
%     3rd party packages
%%%%%%%%%%%%%%%%%%%%%%%%%%%%%%%%%%%%%%
%\usepackage[utf8]{inputenc}
\usepackage[letterpaper, margin=0.5in,headheight=77pt,top=3.5cm]{geometry}
\usepackage{amsmath}
\usepackage{amssymb}
\usepackage{hyperref}
\usepackage[table]{xcolor}
\usepackage{tabu}
\usepackage{fancyhdr}
\usepackage{lastpage}
\usepackage{graphicx}
\usepackage{collcell}
\usepackage{multirow}
\usepackage[skins]{tcolorbox}
\usepackage{fontspec}
    \setmainfont{Lato}

\hypersetup{
  colorlinks   = true, %Colours links instead of using boxes to indicate links
  urlcolor     = black, %Colour for external hyperlinks
  linkcolor    = black %Colour for internal links
  citecolor   = red %Colour for citations
}
\newcolumntype{U}{>{\collectcell\MakeUppercase}l<{\endcollectcell}}

\tcbuselibrary{listings,breakable}


\newtcolorbox{reportSection}[2][]{%
  attach boxed title to top left
  			   = {xshift=10pt, yshift=-8pt},
  colback      = white!5!white,
  colframe     = white!75!black,
  fonttitle    = \bfseries,
  colbacktitle = white,
  title        = \textbf{\large{#2}},
  coltitle	   = black,
  arc = 0mm,
  opacityframe = 0.5,
  boxrule=1pt,
  boxed title style={%
    sharp corners, 
    rounded corners=northwest, 
    colback=white, 
    boxrule=0pt,
    titlerule=0mm},
  enhanced,
}

%--------------------------------------------
% Creating Header Style
%-------------------------------------------


%\renewcommand{\headrule}{\hbox to\headwidth{%
%    \color{gray}\leaders\hrule height \headrulewidth\hfill}}

%--------------------------------------------
% Creating Header with Text, Logo and Page Number
%-------------------------------------------

%-------------- HEADER CONTENT ------------

\pagestyle{fancy}
\fancyhf{}
\lhead{\Large{\textbf{MYCOBACTERIUM TUBERCULOSIS \\ GENOME SEQUENCING REPORT}} \\ %ISO15189 reporting requirements stipulate that you provide "a clear, unambiguous identification of the examination including, where appropriate, the examination procedure"
\normalsize{NOT FOR DIAGNOSTIC USE}} %If you are using this as part of an accredited clinical pipeline, you may remove this text
\chead{}
\rhead{\includegraphics[scale=0.25]{imageFiles/LogoDesign}}
\renewcommand{\headrulewidth}{0pt} % no line in header area
%To accommodate new oversized header
%\setlength{\headheight}{4.5\baselineskip} %this is going to produce an error, but works - HOWEVER, it also makes the footer disappear

%-------------- FOOTER CONTENT ------------
% !!!! WARNING !!!!  OVERLEAF DOESN'T DO A GREAT JOB RUNNING THE \pageref COMMAND THAT WOULD LET US AUTOMATICALLY GET THE LAST PAGE #. SO, THE LAST PAGE IS HARD CODED HERE, THIS WILL CAUSE A PROBLEM IF YOUR REPORT IS LONGER THAN TWO PAGES!!!

\lfoot{Page \thepage\ of 2}
\rfoot{Patient ID: 12345678910 | Date: 2017-01-01 | Location: Someplace }%ISO15189 reporting standards require that patient identification and location occur on every page of the report, along with a 'Page X of Y' designation.



\begin{document}

%---------------------------------------------------------
%  TABLE WITH RELEVANT PATIENT INFORMATION
%--------------------------------------------------------
% Delete this line when including this report in your own pipeline vvvvv
%\noindent \footnotesize{\textbf{See %\url{https://github.com/amcrisan/TB-WGS-MicroReport} for how to automatically fill the contents of this template}}
%\vspace{3mm}
% ^^^^^^^^^^^^^^


\tabulinesep=5pt
\noindent 
\taburulecolor{lightgray}
\noindent \begin{tabu} to \textwidth {|XU||XU|}
\hline
% The table below contains the information that is required by ISO15189 reporting standards
%-------------------- v ADD YOUR TABLE CONTENTS BELOW v --------------------------
Patient Name  & John Doe       & Barcode        &          \\ \hline
Birth Date    & 2000-01-01     & Patient ID          & 12345678910           \\ \hline
Location   & Someplace         & Sample Type       & Sputum          \\ \hline
Sample Source & Pulmonary	   & Sample Date & 2016-12-25      \\ \hline
Sample ID & A12345678	   & Sequenced From & MGIT Cultured Isolate       \\ \hline
Reporting Lab & Lab Name       & Report Date/Time  & 2017-01-01, 15:36   \\ \hline
Requested By  & Requester Name & Requester Contact & requester@email.com \\ \hline
%-------------------- ^ ADD YOUR TABLE CONTENTS ABOVE ^ --------------------------
\end{tabu}
\vspace{5mm}



%---------------------------------------------------------
%  SECTION 1 : SUMMARY OF REPORT CONTENT
%---------------------------------------------------------

\begin{reportSection}{Summary}
\vspace{1mm}

% This summary paragraph presents the most important information in a simple paragraph and uses bolding to indicate the species, predicted drug resistance, and whether or not the isolate clusters with others, suggesting recent transmission. Our research indicated that this brief summary is a feature appreciated by our end-users, particularly busy clinicians
The specimen was positive for \textbf{Mycobacterium tuberculosis}. It is \textbf{resistant to isoniaizd and rifampin}. It belongs to a cluster, suggesting \textbf{recent transmission}.
\end{reportSection}
\vspace{5mm}

%---------------------------------------------------------
%  SECTION 2 : SUMMARY OF SPECIATION RESULTS
%---------------------------------------------------------

\begin{reportSection}{Organism} % Our design study indicated this was the preferred headline word for this section, in which speciation is summarised.
\vspace{1mm}
The specimen was positive for \textit{\textbf{Mycobacterium tuberculosis}}, lineage 2.2.1 (\textbf{East-Asian Beijing}). % Use bold text to highlight the organism name
\end{reportSection}
\vspace{5mm}

%---------------------------------------------------------
%  SECTION 3 : SUMMARY OF DRUG SUSCEPTIBILITY RESULTS
%---------------------------------------------------------


\begin{reportSection}{Drug Susceptibility}

%--------- Text for drug susceptibility and also the summary tick boxes
\vspace{1mm}
\tabulinesep=2pt
\noindent 
\taburulecolor{lightgray}
\noindent \begin{tabu} to \textwidth {X[1.1,l]X}
\multirow{4}{*}{\parbox{7.5cm}{\footnotesize{Resistance is reported when a high-confidence resistance-conferring mutation is detected. \textbf{``No mutation detected'' does not exclude the possibility of resistance}.}}} 
& $\square$ No drug resistance predicted \\
& $\square$ Mono-resistance predicted \\
&  \cellcolor[HTML]{EFEFEF} $\text{\rlap{$\checkmark$}}\square$ Multi-drug resistance predicted \\ % This checkbox summary provides a helpful overview of the predicted resistance phenotype 
& $\square$ Extensive drug resistance predicted \\
\end{tabu}

\vspace{3mm} %% Add a Bit of space between the entities

%--------- Summary of drug susceptibility - table
\tabulinesep=5pt
\noindent 
\taburulecolor{lightgray}
\noindent \begin{tabu} to \textwidth {X[0.5,l]X[0.75,l]X[0.75,l]X[2,l]}
% Our study indicated that including a drug's full name was preferred, as was including the gene in which a resistance-conferring mutation was found. The mutation itself is shown as an amino acid change.
%--------------- v ADD YOUR TABLE CONTENTS BELOW v---------------------
\hline
Drug class                    & Interpretation                                           & Drug                              & 
Resistance Gene \scriptsize{(Amino Acid Mutation)}          \\
\cline{1-4}
                              &                                                     & Ethambutol                        & No mutation detected      \\
                              & \multirow{-2}{*}{Susceptible}                         & Pyrazinimide                      & No mutation detected      \\
                              & \cellcolor[HTML]{EFEFEF}                            & \cellcolor[HTML]{EFEFEF}Isoniazid & \cellcolor[HTML]{EFEFEF}katG (S315T) \\
\multirow{-4}{*}{First Line}  & \multirow{-2}{*}{\cellcolor[HTML]{EFEFEF}Resistant} & \cellcolor[HTML]{EFEFEF}Rifampin  & \cellcolor[HTML]{EFEFEF}rpoB (S531L) \\
\cline{1-4}
                              &                                                     & Streptomycin                      & No mutation detected      \\
                              &                                                     & Ciprofloxacin                     & No mutation detected      \\
                              &                                                     & Ofloxacin                         & No mutation detected      \\
                              &                                                     & Moxifloxacin                       & No mutation detected      \\
                              &                                                     & Amikacin                          & No mutation detected      \\
                              &                                                     & Kanamycin                         & No mutation detected      \\
\multirow{-7}{*}{Second Line} & \multirow{-7}{*}{Susceptible}                         & Capreomycin                       & No mutation detected	\\
\cline{1-4}
%--------------- ^ ADD YOUR TABLE CONTENTS ABOVE ^ ---------------------
\end{tabu}
\end{reportSection}

%---------------------------------------------------------
%  SECTION 4 : CLUSTER DETECTION RESULTS
%---------------------------------------------------------
\pagebreak
\begin{reportSection}{Cluster Detection}
\vspace{3mm}
The current isolate was clustered with previously sequenced isolates, suggesting \textbf{recent transmission}.
\vspace{2mm}

\tabulinesep=5pt
\noindent 
\taburulecolor{lightgray}
\noindent \begin{tabu} to \textwidth {XX}
% Set your own thresholds and choice of language here. You might also wish to populate this section with a pre-written narrative about the cluster your isolate belongs to, if available
%--------------- v ADD YOUR TABLE CONTENTS BELOW v ---------------------
Relatedness  & Number of prior matching isolates  \\ \hline
Closely Related (< 5 mutations apart) & \textbf{2} isolates \\
Related (6 to 30 mutations apart) & \textbf{6} isolates \\ \hline
%--------------- ^ ADD YOUR TABLE CONTENTS ABOVE ^ -----------------
\end{tabu}

\includegraphics[scale=0.65]{./imageFiles/treeExample}

\end{reportSection}
\vspace{1mm}

%---------------------------------------------------------
%  SECTION 6 : ASSAY DETAILS 
%---------------------------------------------------------
\begin{reportSection}{Assay Details}
\vspace{1mm}
\tabulinesep=5pt
\noindent 
\taburulecolor{lightgray}
\noindent \begin{tabu} to \textwidth {|XU||XU|}
\hline
% The table below contains the information that is required by ISO15189 reporting standards
%-------------------- v ADD YOUR TABLE CONTENTS BELOW v --------------------------
Sample ID  & A12345678       & Barcode        &          \\ \hline
Sequencer    & Illumina HiSeq 2500     & Method          & WGS           \\ \hline
Pipeline   & ReSeqTB v.3.2c         & Reference       & H37Rv           \\ \hline
%-------------------- ^ ADD YOUR TABLE CONTENTS ABOVE ^ --------------------------
\end{tabu}
\vspace{1mm}

\end{reportSection}
\vspace{1mm}

%---------------------------------------------------------
%  SECTION 6 : ADDITIONAL COMMENTS
%---------------------------------------------------------
%Comments are for anything funny that was found by the 
\begin{reportSection}{Comments}
% This is where any additional comments might go
\vspace{1mm}

\footnotesize{No additional comments for this report}
\vspace{2mm}

\textcolor{gray}{\footnotesize{\textit{\textbf{Standard Disclaimer:} Low frequency hetero-resistance below the limit of detection by sequencing may affect typing results.  The interpretation provided is based on the current understanding of genotype-phenotype relationships.}}}

\end{reportSection}
\vspace{1mm}


%---------------------------------------------------------
%  SECTION 7 : DISCLAIMER
%---------------------------------------------------------
% Note: Disclaimer is a general note that will appear on all reports. In contrast to Comments, which is for the the result of the individual patient. 
%\begin{reportSection}{Disclaimer}
% This is where any additional comments might go
%\vspace{1mm}

%\footnotesize{Low frequency hetero-resistance below the limit of detection by sequencing may affect typing results.  The interpretation provided is based on the current understanding of genotype-phenotype relationships.}
%\end{reportSection}
%\vspace{1mm}

%---------------------------------------------------------
%  SECTION 8 : AUTHORIZATION
%---------------------------------------------------------
\begin{reportSection}{Authorised}
% This section is required by ISO15189 standards
\vspace{1mm}
\tabulinesep=5pt
\noindent 
\taburulecolor{lightgray}
\noindent \begin{tabu} to \textwidth {|XU||XU|}
\hline
%--------------- v ADD YOUR TABLE CONTENTS BELOW v ---------------------
Signature  &   & Name  &   \\ \hline
Position   &   & Date    &        \\ \hline
%--------------- ^ ADD YOUR TABLE CONTENTS ABOVE ^ ---------------------
\end{tabu}

\end{reportSection}



\end{document}
